%% =======================================================
%%  Manuscript prepared for submission to
%%  Expert Systems with Applications (Elsevier)
%%
%%  CURRENT STATUS:
%%   Numeric and factual placeholders available from completed experiments have
%%   been filled. MotionAGFormer-XS and MotionBERT-Lite (81-frame) are complete.
%% =======================================================

\documentclass[preprint,12pt]{elsarticle}

\usepackage{amsmath}
\usepackage{amssymb}
\usepackage{graphicx}
\usepackage{booktabs}
\usepackage{multirow}
\usepackage{hyperref}
\usepackage{float}
\usepackage{microtype}
\usepackage{lineno}
\usepackage{xcolor}
\usepackage{diagbox}

\journal{Expert Systems with Applications}

\begin{document}

\begin{frontmatter}

%% =======================================================
%%  TITLE
%% =======================================================
\title{An intelligent vision-based decision-support system for
Parkinson's gait severity assessment using COM-centred
spatio-temporal modeling from a single RGB camera}

%% =======================================================
%%  AUTHORS
%% =======================================================
\author[cnu]{Bokyung Kim}
\author[cnu]{Hieyong Jeong\corref{cor1}}
\ead{h.jeong@jnu.ac.kr}
\author[cnu]{Md Ilias Bappi}
\author[cnu]{Kyungbaek Kim}

\cortext[cor1]{Corresponding author}

\address[cnu]{Chonnam National University, Gwangju, Republic of Korea}

%% =======================================================
%%  ABSTRACT
%% =======================================================
\begin{abstract}
Parkinson's disease (PD) is a progressive neurodegenerative disorder for
which gait impairment serves as a key clinical biomarker. In current
practice, gait severity is rated on the MDS-UPDRS Part~III scale through
expert visual judgement, which is inherently subjective, episodic, and
limited to clinical settings. An automated and deployable assessment tool
is therefore needed to support clinicians in objective decision making.

This paper proposes an intelligent decision-support system estimating
the MDS-UPDRS Part~III gait score directly from a single frontal-view RGB
video. The system integrates four key components: (i) a Centre-of-Mass (COM)
centred spatial normalisation scheme that removes global subject-position
variation and is further analysed with an optional bone-length scaling variant
for camera-distance robustness;
(ii) a multi-channel kinematic feature representation; (iii) a Graph
Convolutional Spatio-Temporal Encoder with learnable adjacency; and (iv) a
selective-prediction module coupled with bounded regression that enforces the
observable clinical range and defers borderline cases to human review.

We develop and evaluate the system using a dual-cohort strategy. A
110-patient subset (6{,}066 sequences) from the multi-site CARE-PD benchmark
serves as the primary training and cross-validation domain. A prospectively
collected independent clinical cohort (CNUH, $N=21$) is used strictly for
cross-site zero-shot transfer, simulating deployment in an unseen environment.
All competing models are evaluated under an identical input representation and
data split, and we report both Mean Absolute Error (MAE) and Macro-F1 to
expose, rather than mask, the strong class imbalance of the ordinal scale.
Under subject-level GroupKFold cross-validation the system attains
MAE\,=\,0.358, RMSE\,=\,0.564, and MedAE\,=\,0.147, and performs competitively
with state-of-the-art motion encoders in both MAE and Macro-F1, while providing
bounded-output safety and an explicit clinical abstention rule. A scale-normalised
robustness variant additionally demonstrates invariance to controlled
coordinate scaling under a global-scaling approximation.
Zero-shot transfer (CARE-PD\,$\to$\,CNUH) demonstrates stable cross-site
behaviour, and an extreme low-data stress test confirms that bounded
regression prevents the catastrophic extrapolation observed in unbounded
baselines. By additionally reporting inference latency and a quantitative
risk--coverage analysis of the abstention rule, the proposed pipeline
establishes a practical foundation for privacy-preserving, edge-deployable PD
gait monitoring.
\end{abstract}

%% =======================================================
%%  KEYWORDS
%% =======================================================
\begin{keyword}
Intelligent decision-support system \sep
Parkinson's disease \sep
Vision-based gait analysis \sep
Graph convolutional network \sep
Selective prediction \sep
MDS-UPDRS
\end{keyword}

\end{frontmatter}

\linenumbers

%% =======================================================
%%  HIGHLIGHTS  (paste into the separate Highlights file; each <=85 chars)
%% =======================================================
% * Vision-based decision-support system for MDS-UPDRS gait scoring from RGB video.
% * COM-centred normalisation reduces global subject-position variation.
% * Graph Convolutional Spatio-Temporal Encoder with bounded-output safety head.
% * Dual-cohort study: CARE-PD (training) and CNUH (zero-shot cross-site transfer).
% * Selective-prediction rule and edge latency reported for clinical deployment.

%% =======================================================
%%  1. INTRODUCTION
%% =======================================================
\section{Introduction}
\label{sec:intro}

\subsection{Background and clinical motivation}
Parkinson's disease (PD) is characterised by progressive motor symptoms
including bradykinesia, rigidity, and impaired gait \cite{goetz2008mds}.
Gait severity is clinically rated using Part~III of the Movement Disorder
Society's Unified Parkinson's Disease Rating Scale (MDS-UPDRS)
\cite{martinez2011mds}. This rating relies on expert visual judgement,
exhibits non-trivial inter-rater variability \cite{richards1994inter}, and
is obtained only during scheduled visits, thereby missing the intraday
symptom fluctuations driven by medication cycles \cite{zhang2024wearable}.
These limitations motivate \emph{intelligent decision-support systems} that
can deliver objective, reproducible, and continuous quantification of gait
severity to assist clinicians.

\subsection{Why an RGB-only expert system?}
Marker-based optical motion-capture systems offer millimetre-level accuracy
but require dedicated laboratories \cite{dibiase2020gait}, while wearable
inertial sensors demand strict placement, charging, and skin-contact
adherence \cite{zhang2024wearable}. By contrast, an expert system that
estimates MDS-UPDRS gait scores directly from a single RGB camera meets
several practical deployment desiderata: zero hardware burden on the patient,
compatibility with existing clinic and consumer cameras, the potential for
privacy-preserving edge deployment when only skeletal keypoints are stored,
and scalability across sites \cite{sun2024digital}.

\subsection{Open problems and our contributions}
Despite progress in skeleton-based motion modelling
\cite{shi2019twoStream,zhu2023motionbert}, deploying PD gait assessment as a
reliable clinical system requires addressing three concrete gaps:

\begin{enumerate}
    \item \textbf{Robustness to environmental variation.} Predictions from a
    single RGB video are sensitive to camera-to-subject distance (apparent
    scale) and horizontal subject position; few algorithmic studies formulate
    and empirically quantify an explicit invariance mechanism against these
    spatial confounders \cite{gennarelli2025skelmamba}.
    \item \textbf{Evaluation bias in ordinal assessment.} Vision-based
    regression formulations frequently report Mean Absolute Error (MAE) alone,
    which can mask regression-to-the-mean behaviour on minority severe cases;
    the classification benchmarks of this field instead report Macro-F1
    \cite{adeli2025carepd}.
    \item \textbf{Clinical decision support versus black-box prediction.}
    A fielded expert system must provide deployment transparency, operational
    rules such as bounded clinical ranges, abstention mechanisms for uncertain
    cases, and measured computational cost \cite{sun2024digital}.
\end{enumerate}

To address these gaps, we propose an intelligent decision-support system with
the following contributions:
\begin{itemize}
    \item We formalise a \textbf{COM-centred spatial normalisation scheme}
    that removes global subject-position variation, and we separately evaluate
    an optional median bone-length scaling variant for controlled
    camera-distance robustness (Section~\ref{sec:robust_exp}).
    \item We process multi-channel kinematic features with a \textbf{Graph
    Convolutional Spatio-Temporal Encoder} featuring learnable adjacency
    \cite{shi2019twoStream}, discovering non-anatomical inter-joint
    dependencies indicative of pathological gait compensations.
    \item We evaluate rigorously with both \textbf{MAE and Macro-F1}
    (reported as F1$_{0\text{--}3}$ and F1$_{0\text{--}2}$ to match the CARE-PD
    benchmark), transparently handling class imbalance, and we run every
    competing model under an \textbf{identical input representation and data
    split} so that comparisons reflect architecture rather than preprocessing.
    \item We explicitly evaluate \textbf{zero-shot cross-site generalisation}
    (CARE-PD\,$\to$\,independent CNUH cohort) and show that the
    \textbf{bounded-regression} safety head prevents the catastrophic
    extrapolation seen in unbounded baselines under domain shift.
    \item We formalise a \textbf{selective-prediction (abstention) rule},
    quantify it with a risk--coverage analysis against a random-abstention
    baseline, and report practical \textbf{edge-inference latency}
    (Section~\ref{sec:discussion}).
\end{itemize}

%% =======================================================
%%  2. RELATED WORK
%% =======================================================
\section{Related Work}
\label{sec:related}

\subsection{Vision-based gait analysis for PD}
The maturity of markerless pose estimation \cite{cao2017openpose,
bazarevsky2020blazepose} has enabled non-contact PD gait analysis from RGB
video. Lu et al. \cite{lu2020vision,lu2021quantifying} pioneered vision-based
MDS-UPDRS gait scoring from 3D pose sequences. The CARE-PD benchmark
\cite{adeli2025carepd} consolidates 363 participants across nine cohorts from
eight clinical sites, distributed as harmonised SMPL meshes, and defines
standardised generalisation protocols (within-dataset, cross-dataset,
leave-one-dataset-out, and in-domain adaptation), enabling large-scale
validation of motion encoders \cite{gennarelli2025skelmamba,
tian2024crossspatiotemporal,adeli2024benchmarking}. We utilise a curated
110-participant, fully annotated subset of CARE-PD as our primary training
domain. A recurring limitation across this literature is the unquantified
sensitivity of predictions to camera scale and subject distance.

\subsection{Skeleton-based deep models and expert systems}
Skeleton modelling has evolved from fixed-graph convolutions
(ST-GCN \cite{yan2018spatial}) to adaptive graphs
(2s-AGCN \cite{shi2019twoStream}) and Transformer-based encoders
(MotionBERT \cite{zhu2023motionbert},
MotionAGFormer \cite{mehraban2024motionagformer}). From an expert-systems
perspective, however, raw accuracy is insufficient: a fielded system must also
exhibit domain-specific robustness, interpretable failure modes, an explicit
uncertainty/abstention policy, and edge-level efficiency. Our work bridges
this gap by embedding a biomechanically motivated normalisation and
deployment-centric decision rules into a modern spatio-temporal architecture.

%% =======================================================
%%  3. METHODOLOGY
%% =======================================================
\section{Proposed System}
\label{sec:method}

The proposed system takes a single frontal-view RGB gait video and outputs a
continuous MDS-UPDRS Part~III (Item~3.10) gait estimate. It comprises four
stages: COM-centred normalisation, multi-channel feature construction,
spatio-temporal encoding, and a bounded-regression head.

\subsection{Stage 1: Pose extraction and COM-centred normalisation}
\label{sec:stage1}

\paragraph{Pose estimation.}
We apply the MediaPipe Pose estimator \cite{bazarevsky2020blazepose} to each
frame, obtaining per-joint image-plane coordinates and a model-predicted
relative depth. We treat the result as a 2.5D coordinate
$\mathbf{p}^{\text{raw}}_{t,j}\in\mathbb{R}^{3}$, used as structured input
rather than as a calibrated 3D measurement, with $T=390$ frames and $J=17$
COCO-format joints.
For CARE-PD, we used the released SMPL-derived H36M-compatible skeleton
representation and converted it to the same 17-joint layout used throughout the
experiments. For CNUH, MediaPipe landmarks were mapped to the same
H36M-compatible 17-joint format, yielding a common input representation across
cohorts.

\paragraph{COM-centred normalisation.}
To remove the subject's global image-plane position, we approximate the body
centre (COM proxy) as the midpoint of the left and right hips, which is
biomechanically co-located with the trunk's centre of mass during quiet
bipedal locomotion \cite{winter2009biomechanics}:
\begin{equation}
\mathbf{C}_t = \tfrac{1}{2}\!\left(
\mathbf{p}^{\text{raw}}_{t,\text{L-hip}}+\mathbf{p}^{\text{raw}}_{t,\text{R-hip}}\right),
\qquad
\mathbf{p}^{\text{rel}}_{t,j}=\mathbf{p}^{\text{raw}}_{t,j}-\mathbf{C}_t .
\label{eq:trans}
\end{equation}
The main predictive model uses $\mathbf{p}^{\text{rel}}_{t,j}$ as its coordinate
basis. To test robustness to apparent-scale variation induced by camera
distance, we additionally evaluate an optional median bone-length normalised
variant:
\begin{equation}
\mathbf{p}^{\text{scale}}_{t,j}=\frac{\mathbf{p}^{\text{rel}}_{t,j}}{L_{\text{ref}}},
\label{eq:norm}
\end{equation}
where $L_{\text{ref}}$ is the median torso bone length computed over the
sequence. Under a global multiplicative model of distance change,
$\mathbf{p}^{\text{scale}}$ is invariant to coordinate scaling by construction;
Section~\ref{sec:robust_exp} verifies this empirically and discusses its limits
under true perspective distortion. This scale-normalised variant is used only
for the robustness analysis, while the main performance tables report the
best-performing COM-centred model.

\subsection{Stage 2: Multi-channel kinematic feature construction}
\label{sec:stage2}
From $\mathbf{P}_{\text{rel}}$ we construct a per-(frame, joint) feature
vector by concatenating clinically motivated channels---joint positions
($f_{\text{pos}}$), frame-difference velocities ($f_{\text{vel}}$,
sensitive to bradykinesia), sequence-level motion amplitude and variability
($f_{\text{amp}},f_{\text{var}}$), and joint angles
($f_{\text{ang}}$, encoding range of motion):
\begin{equation}
\mathbf{x}_{t,j}=\mathrm{Concat}\!\left[
f_{\text{pos}},f_{\text{vel}},f_{\text{amp}},f_{\text{var}},f_{\text{ang}}\right]\in\mathbb{R}^{C}.
\end{equation}

\subsection{Stage 3: Graph Convolutional Spatio-Temporal Encoder}
\label{sec:stage3}
Unlike ST-GCN \cite{yan2018spatial}, which relies on a fixed anatomical graph,
we use a \emph{learnable} adjacency matrix \cite{shi2019twoStream} to discover
task-relevant, non-anatomical inter-joint dependencies characteristic of
pathological gait compensation. Two GraphConv layers (with LayerNorm and ReLU)
expand the per-joint dimension ($C\!\to\!64\!\to\!128$). A multi-head
self-attention module over the $J=17$ joint tokens captures global spatial
relations, and a Temporal Transformer encoder along the time axis models
frame-to-frame dynamics including stride periodicity. The sequence is collapsed
by temporal average pooling to $\mathbf{h}\in\mathbb{R}^{J\times128}$.

\subsection{Stage 4: Bounded regression head}
\label{sec:stage4}
To enforce clinical validity, the pooled representation is mapped by an MLP to a
sigmoid-scaled bounded output:
\begin{equation}
\hat{y}=3\cdot\sigma\!\left(\mathrm{MLP}(\mathbf{h})\right)\in(0,3).
\label{eq:bounded}
\end{equation}
Because score~4 denotes inability to walk independently---which precludes
frontal-view ambulatory video---the range $[0,3]$ covers all feasible cases in
both cohorts. The network is trained by minimising the MAE between $\hat{y}$ and
the ground-truth score $y$.

\subsection{Implementation and reproducibility}
\label{sec:repro}
Table~\ref{tab:hparams} reports all hyperparameters. Source code and
configuration files are released at
\url{https://github.com/bokyung08/parkinson_COM}.

\begin{table}[htpb]
\centering
\caption{Implementation and training hyperparameters of the proposed system.}
\label{tab:hparams}
\footnotesize
\setlength{\tabcolsep}{4pt}
\renewcommand{\arraystretch}{0.95}
\resizebox{0.95\linewidth}{!}{%
\begin{tabular}{ll}
\toprule
\textbf{Component} & \textbf{Setting} \\
\midrule
Pose estimator               & MediaPipe Pose (BlazePose Heavy) \\
Joints / Sequence length     & $J=17$ (COCO) / $T=390$ (interpolation) \\
Normalisation                & COM translation (bone-length scaling only in robustness variant) \\
GraphConv hidden dim         & $C \to 64 \to 128$ (learnable adjacency) \\
Joint-attention heads        & 4 \\
Temporal Transformer         & 2 layers, 4 heads \\
MLP hidden dim / Dropout     & 256 / 0.1 \\
Output activation            & $3\cdot\sigma(\cdot)$ (bounded regression) \\
Total parameters             & 158{,}594 \\
\midrule
Optimiser / LR / Batch       & Adam / $1\times10^{-4}$ / 16 \\
Epochs / Loss                & 100 / MAE \\
Framework / Hardware         & PyTorch 2.0 / NVIDIA RTX 3080 \\
\bottomrule
\end{tabular}
}
\end{table}

%% =======================================================
%%  4. EXPERIMENTAL SETUP
%% =======================================================
\section{Experimental Setup}
\label{sec:setup}

\subsection{Datasets and cohort roles}
\label{sec:datasets}
\paragraph{Training benchmark (CARE-PD subset).}
For model development and cross-validation we use a fully annotated subset of
the CARE-PD benchmark \cite{adeli2025carepd} comprising 110 participants and
6{,}066 sequences drawn from four cohorts (3DGait, BMCLab, PD-GaM, T-SDU-PD).

\paragraph{Zero-shot transfer cohort (CNUH).}
To evaluate cross-site transferability, frontal-view gait videos were
prospectively collected from $N=21$ PD patients at Chonnam National University
Hospital (IRB: CNUH-2025-203; written informed consent obtained). This cohort
was held out entirely from training and is used only as an unseen target
domain.

\subsection{Evaluation protocol and metrics}
\label{sec:metrics}
We adopt subject-level GroupKFold 5-fold cross-validation on the CARE-PD subset,
ensuring that all sequences of a subject lie in a single fold. We report MAE and
RMSE in score units, and Macro-F1 obtained by rounding the continuous output to
the nearest integer (clipped to $[0,3]$). Following the CARE-PD benchmark, we
report two Macro-F1 variants: F1$_{0\text{--}3}$ over all classes and
F1$_{0\text{--}2}$ excluding the rare score-3 class (the metric only; all
score-3 samples remain in training). We note that rounding a regressor is not
identical to a native ordinal classifier, so Macro-F1 comparisons with
classification benchmarks are indicative rather than exact. For significance
testing we use the Wilcoxon signed-rank test. To avoid inflating significance by
treating multiple sequences from the same subject as independent samples,
absolute errors were also aggregated at the subject level before testing,
yielding $p=0.107$ for Ours versus Lu et al. and
$p=9.09\times10^{-4}$ for Ours versus ST-GCN.

\subsection{Fair-comparison protocol}
\label{sec:fairness}
All competing models---classical regressors, ST-GCN \cite{yan2018spatial},
Lu et al.\ \cite{lu2020vision}, MotionBERT \cite{zhu2023motionbert}, and
MotionAGFormer \cite{mehraban2024motionagformer}---are trained and evaluated on
the \emph{identical} MediaPipe 17-joint input, the COM-centred
representation of Section~\ref{sec:stage1}, and the same subject-level
GroupKFold split. Transformer encoders were re-trained on this input with their
native temporal resampling applied to reach $T=390$; publicly released
checkpoints were used for initialisation where available. Consequently, the
spatial-robustness property of Section~\ref{sec:robust_exp} is a property of
the \emph{shared preprocessing pipeline}, and the comparison in
Section~\ref{sec:comparison} isolates architectural differences alone.
For MotionBERT-Lite, we use an 81-frame uniformly sampled temporal adapter to
reduce the quadratic cost of temporal self-attention while preserving the same
subject-level split, training epochs, and 17-joint representation.

%% =======================================================
%%  5. EXPERIMENTAL RESULTS
%% =======================================================
\section{Experimental Results}
\label{sec:results}

\subsection{Quantitative robustness to spatial perturbations}
\label{sec:robust_exp}
To validate the optional scale-normalisation step in Eq.~\eqref{eq:norm}, we
apply synthetic scale (distance) and horizontal-translation perturbations to the
input coordinates at inference time (Table~\ref{tab:robustness}).

\begin{table}[htpb]
\centering
\caption{Robustness to affine scale (distance) and translation perturbations.
$\Delta$MAE is the relative change from each model's own unperturbed baseline.}
\label{tab:robustness}
\footnotesize
\setlength{\tabcolsep}{4pt}
\renewcommand{\arraystretch}{0.95}
\resizebox{0.95\linewidth}{!}{%
\begin{tabular}{lcc}
\toprule
\textbf{Perturbation} & \textbf{COM only ($\Delta$MAE)} &
\textbf{COM + bone-norm variant ($\Delta$MAE)} \\
\midrule
Original (baseline MAE)        & \textbf{0.358} & 0.366 \\
Scale $s=0.70$ (far)           & $+35.9\%$ & \textbf{$+0.0\%$} \\
Scale $s=1.30$ (near)          & $+46.5\%$ & \textbf{$+0.0\%$} \\
Shift $\Delta x=\pm0.20$       & $+0.0\%$  & \textbf{$+0.0\%$} \\
\bottomrule
\end{tabular}
}
\end{table}

Without bone-length scaling, distance-induced scale perturbations degrade MAE by
up to $46.5\%$. Adding the scaling step eliminates this synthetic shift
($\Delta$MAE $=0\%$). This invariance holds \emph{by construction} under a global
multiplicative model of distance change, and we do not claim robustness to true
perspective distortion arising from extreme camera angles, which we treat as a
limitation (Section~\ref{sec:failure}). Because the scale-normalised variant has
a small accuracy cost (0.358\,$\to$\,0.366 MAE), the main comparison tables use
the best-performing COM-centred model, while the scale-normalised variant is
reported as a robustness operating point.

\subsection{Ablation studies}
\label{sec:ablation}
Table~\ref{tab:ablation_feat} reports the input-feature ablation and
Table~\ref{tab:ablation_arch} the architecture ablation, both on the combined CNUH+CARE-PD cohort under the GroupKFold protocol.

\begin{table}[htpb]
\centering
\caption{Input-feature ablation on the combined CNUH+CARE-PD cohort
(GroupKFold 5-fold, $N=6{,}087$).
Lower MAE/RMSE and higher F1 are better.}
\label{tab:ablation_feat}
\begin{tabular}{clccc}
\toprule
\textbf{Config.} & \textbf{Feature set} & \textbf{MAE} & \textbf{RMSE} &
\textbf{F1$_{0\text{--}3}$} \\
\midrule
A & COM coordinates only            & 0.374 & 0.577 & 0.617 \\
B & + velocity                      & 0.432 & 0.629 & 0.508 \\
C & + amplitude / variability       & 0.376 & 0.549 & 0.546 \\
D & \textbf{+ joint angle (Proposed)} & \textbf{0.358} & 0.564 & \textbf{0.661} \\
\bottomrule
\end{tabular}
\end{table}

\begin{table}[htpb]
\centering
\caption{Architecture ablation on the combined CNUH+CARE-PD cohort
(Configuration~D features, GroupKFold 5-fold).}
\label{tab:ablation_arch}
\begin{tabular}{lccc}
\toprule
\textbf{Model} & \textbf{MAE} & \textbf{RMSE} & \textbf{F1$_{0\text{--}3}$} \\
\midrule
MLP only (mean pooling)            & 0.554 & 0.653 & 0.381 \\
GraphConv + MLP                    & 0.450 & 0.580 & 0.510 \\
GraphConv + Joint Attention + MLP  & 0.414 & 0.564 & 0.507 \\
\textbf{Full model (Ours)}         & \textbf{0.358} & 0.564 & \textbf{0.661} \\
\bottomrule
\end{tabular}
\end{table}

\subsection{Comparison with state-of-the-art motion encoders}
\label{sec:comparison}
Table~\ref{tab:comparison} compares the proposed system against classical,
graph-based, and Transformer-based baselines, all under the fair-comparison
protocol of Section~\ref{sec:fairness}.

\begin{table}[htpb]
\centering
\caption{Comparison on the combined CNUH+CARE-PD cohort (17-joint, COM-centred input,
subject-level GroupKFold 5-fold, $N=6{,}087$). All models share the identical
input and split; the comparison therefore reflects architecture alone.}
\label{tab:comparison}
\begin{tabular}{lccccc}
\toprule
\textbf{Model} & \textbf{\#Params} & \textbf{MAE} & \textbf{RMSE} &
\textbf{F1$_{0\text{--}3}$} & \textbf{F1$_{0\text{--}2}$} \\
\midrule
SVR (best classical)              & ---     & 0.492 & 0.639 & 0.485 & 0.597 \\
Lu et al.\ \cite{lu2020vision}    & 147{,}908 & 0.404 & 0.543 & 0.666 & 0.673 \\
ST-GCN \cite{yan2018spatial}      & 252{,}097 & 0.443 & 0.623 & 0.476 & 0.628 \\
MotionBERT-Lite (81-frame) \cite{zhu2023motionbert}        & 10.8M & 0.442 & 0.625 & 0.457 & 0.613 \\
MotionAGFormer-XS \cite{mehraban2024motionagformer} & 2.3M & 0.405 & 0.638 & 0.561 & 0.636 \\
\midrule
\textbf{Proposed System (Ours)}   & 158{,}594 & \textbf{0.358} & 0.564 & 0.661 & \textbf{0.691} \\
\bottomrule
\end{tabular}
\end{table}

With only 158{,}594 parameters, the proposed system performs competitively with
substantially larger Transformer encoders in both MAE and Macro-F1 while adding
bounded-output safety. Because all baselines receive the identical COM-centred
input, any accuracy difference is attributable to architecture rather than
preprocessing.
Among the completed baselines, Lu et al. attains the highest rounded-score
F1$_{0\text{--}3}$ (0.666) and the lowest RMSE, whereas the proposed system
achieves the lowest MAE (0.358), strong F1$_{0\text{--}3}$ (0.661), the highest
F1$_{0\text{--}2}$ among completed models (0.691), bounded-output safety, and
substantially lower parameter count than large pretrained Transformer encoders.
MotionAGFormer-XS achieves a similar MAE to Lu et al. (0.405) and the lowest
MedAE among completed external encoders (0.095; reported in the supplementary
result table), but its RMSE and rounded-score F1 are lower than those of the
proposed system. MotionBERT-Lite with the 81-frame adapter is computationally
tractable but remains less accurate than the proposed system and Lu et al.

\subsection{Cross-site zero-shot transfer and bounded safety}
\label{sec:crossds}

\paragraph{Primary zero-shot transfer (CARE-PD\,$\to$\,CNUH).}
To assess real-world deployability, the model is trained on the CARE-PD subset
($N=110$) and evaluated, without any fine-tuning, on the independent CNUH cohort
($N=21$) (Table~\ref{tab:zero_primary}).

\begin{table}[htpb]
\centering
\caption{Primary zero-shot transfer (CARE-PD\,$\to$\,CNUH). No fine-tuning or
domain adaptation was applied.}
\label{tab:zero_primary}
\begin{tabular}{lccc}
\toprule
\textbf{Model} & \textbf{MAE} & \textbf{RMSE} & \textbf{F1$_{0\text{--}3}$} \\
\midrule
ST-GCN \cite{yan2018spatial}      & 1.203 & 1.385 & 0.100 \\
Lu et al.\ \cite{lu2020vision}    & 0.865 & 1.027 & 0.155 \\
Proposed (Ours)          & 0.910 & 1.034 & 0.178 \\
\bottomrule
\end{tabular}
\end{table}

\paragraph{Extreme low-data stress test (CNUH\,$\to$\,CARE-PD).}
To isolate the effect of the bounded-regression head, we reverse the direction:
training on only 21 CNUH subjects and extrapolating to 6{,}066 CARE-PD sequences
(Table~\ref{tab:zero_stress}). This is a deliberate stress test of numerical
stability under severe data scarcity, not a deployment scenario.

\begin{table}[htpb]
\centering
\caption{Extreme low-data stress test (CNUH\,$\to$\,CARE-PD), isolating the
bounded-output safety mechanism.}
\label{tab:zero_stress}
\begin{tabular}{lcc}
\toprule
\textbf{Model variant} & \textbf{MAE} & \textbf{RMSE} \\
\midrule
ST-GCN (unbounded)            & 8.346 & 9.737 \\
ST-GCN (post-hoc clipped to $[0,3]$)       & 2.211 & 2.347 \\
\textbf{Proposed (bounded)}   & \textbf{0.747} & \textbf{0.882} \\
\bottomrule
\end{tabular}
\end{table}

The unbounded ST-GCN suffers catastrophic failure (MAE 8.346) from extreme
extrapolation on domain-shifted data. Post-hoc clipping of its predictions to
the valid range $[0,3]$ prevents this explosion, confirming that bounded output
is essential for numerical safety; the proposed architecture additionally
retains the lowest error, showing that bounding and architecture contribute
separately.
Post-hoc clipping prevents the numerical explosion of unbounded ST-GCN but
remains substantially worse than the proposed bounded model; thus, output
bounding improves safety, while the proposed architecture is still needed for
low error.

\subsection{Clinical performance and transparency}
\label{sec:perclass}
Table~\ref{tab:perclass} stratifies error by ground-truth score, and
Table~\ref{tab:lodo} reports the CARE-PD leave-one-dataset-out (LODO)
protocol, a cohort-level generalisation test.

\begin{table}[htpb]
\centering
\caption{Per-class error and per-class F1 on the combined CNUH+CARE-PD cohort
(GroupKFold 5-fold).}
\label{tab:perclass}
\begin{tabular}{crccc}
\toprule
\textbf{Score} & \textbf{$N$ (seq.)} & \textbf{MAE} & \textbf{RMSE} &
\textbf{Class F1} \\
\midrule
0 & 2{,}615 & 0.225 & 0.395 & 0.793 \\
1 & 2{,}183 & 0.342 & 0.482 & 0.662 \\
2 & 1{,}244 & 0.649 & 0.889 & 0.612 \\
3 &    45  & 0.738 & 0.930 & 0.576 \\
\bottomrule
\end{tabular}
\end{table}
The row-normalised confusion matrix and calibration-style curve are provided as
supplementary transparency figures.

\begin{table}[htpb]
\centering
\caption{CARE-PD leave-one-dataset-out (LODO) results. Each row holds out one
source cohort; Overall aggregates all held-out sequences.}
\label{tab:lodo}
\begin{tabular}{lccc}
\toprule
\textbf{Held-out cohort} & \textbf{MAE} & \textbf{RMSE} &
\textbf{F1$_{0\text{--}3}$} \\
\midrule
3DGait    & 0.775 & 0.947 & 0.233 \\
BMCLab    & 0.663 & 0.844 & 0.340 \\
PD-GaM    & 0.495 & 0.724 & 0.414 \\
T-SDU-PD  & 0.692 & 0.836 & 0.305 \\
\midrule
\textbf{Overall} & \textbf{0.620} & \textbf{0.813} & \textbf{0.358} \\
\bottomrule
\end{tabular}
\end{table}

%% =======================================================
%%  6. DISCUSSION: SYSTEM DEPLOYMENT
%% =======================================================
\section{Discussion: A Clinical Decision-Support System}
\label{sec:discussion}

To qualify as an expert system rather than a bare predictor, the model must
manage its own uncertainty and report concrete deployment cost. We address both.

\subsection{Selective prediction: a risk--coverage analysis}
\label{sec:selective}
A deployable clinical system should defer unreliable predictions to a human
expert. We equip the system with a \emph{post-hoc abstention rule} requiring no
retraining. Because the rounded decision boundaries of the ordinal scale lie at
the half-integers $\{0.5,1.5,2.5\}$, the distance of a continuous prediction
$\hat{y}$ to its nearest boundary is a natural confidence proxy:
\begin{equation}
u(\hat{y}) = \min_{b\in\{0.5,1.5,2.5\}} \lvert \hat{y}-b\rvert .
\label{eq:uncertainty}
\end{equation}
A small $u$ indicates a boundary-proximal prediction most likely to be misrounded;
a large $u$ indicates a prediction comfortably inside a class interval. We
\emph{abstain} (flag for clinician review) whenever $u(\hat{y})<\tau$ and
auto-score the remainder. Sweeping $\tau$ traces a \emph{risk--coverage curve},
where \emph{coverage} is the fraction of cases auto-scored and \emph{selective
risk} is the MAE on the retained subset.

Discarding hard cases trivially lowers retained error, so to demonstrate that
the gain reflects an \emph{informative} uncertainty signal we compare against a
\emph{random-abstention} baseline that flags the same fraction of cases
uniformly at random; an informative rule must lie strictly below this baseline
at equal coverage. We summarise the full curve by the Area Under the
Risk--Coverage curve (AURC; lower is better). We additionally report a fixed,
clinically interpretable operating band $\hat{y}\in[1.3,1.7]$, which targets the
contentious $1\!\leftrightarrow\!2$ boundary where inter-rater disagreement is
highest \cite{richards1994inter}.

\begin{table}[htpb]
\centering
\caption{Selective prediction (risk--coverage) on the combined CNUH+CARE-PD
cohort. At each
target coverage, the threshold $\tau$ flags the most boundary-proximal cases via
Eq.~\eqref{eq:uncertainty}. We report MAE and F1$_{0\text{--}3}$ on the retained
(auto-scored) subset, and the MAE of a random-abstention baseline at matched
coverage. The final row is the fixed clinical band $[1.3,1.7]$.}
\label{tab:selective}
\scriptsize
\setlength{\tabcolsep}{3pt}
\renewcommand{\arraystretch}{0.92}
\resizebox{\linewidth}{!}{%
\begin{tabular}{lccccc}
\toprule
\textbf{Coverage} & \textbf{$\tau$} & \textbf{Flagged (\%)} &
\textbf{Retained MAE} & \textbf{Retained F1$_{0\text{--}3}$} &
\textbf{Random-abst.\ MAE} \\
\midrule
100\% (no abstention) & ---          & 0.0           & 0.358        & 0.661 & 0.358 \\
90\%                  & 0.11  & 10.0          & 0.325 & 0.687 & 0.358 \\
80\%                  & 0.22  & 20.0          & 0.288 & 0.692 & 0.358 \\
70\%                  & 0.30  & 30.0          & 0.260 & 0.701 & 0.357 \\
\midrule
Clinical band $[1.3,1.7]$ & ---      & 6.0    & 0.345 & 0.675 & 0.358 \\
\midrule
\multicolumn{6}{l}{AURC (Ours): 0.278 \quad vs.\ random abstention:
0.358 \quad (lower is better)} \\
\bottomrule
\end{tabular}
}
\end{table}

As Table~\ref{tab:selective} shows, retained MAE decreases from 0.358 at full
coverage to 0.260 at 70\% coverage (a 27.2\% relative reduction), while the
random-abstention baseline at the same coverage remains essentially flat
(0.357--0.358). The gap confirms that the boundary-distance score of
Eq.~\eqref{eq:uncertainty} is an informative uncertainty proxy rather than a
trivial discarding of hard samples. The fixed clinical band $[1.3,1.7]$ flags
only 6.0\% of cases yet reduces retained MAE to 0.345, giving a simple,
auditable human-in-the-loop operating point; the AURC of 0.278 versus 0.358 for
random abstention quantifies the benefit over the entire curve.

\subsection{Edge deployability}
\label{sec:deploy}
With 158{,}594 parameters, the model requires 4.66\,ms per sequence on CUDA GPU
and 5.79\,ms per sequence on CPU in a forward-pass benchmark with batch size 32,
enabling local inference without cloud infrastructure. Once skeletal keypoints are extracted, the raw RGB stream is
discarded, so only low-dimensional joint coordinates are stored or transmitted,
substantially reducing re-identification risk relative to pixel-based pipelines.

\subsection{Failure modes and limitations}
\label{sec:failure}
\begin{itemize}
    \item \emph{Class imbalance.} Score-3 sequences are severely
    underrepresented ($\sim$0.7\%), limiting discriminative capacity for the most
    severe cases; class-weighted losses or targeted augmentation are natural
    extensions.
    \item \emph{Perspective distortion.} COM-centred normalisation removes
    global translation, and the optional bone-length scaling variant neutralises
    global 2D scale perturbations \emph{by construction}, but neither models
    full 3D perspective distortion from extreme camera angles; robustness under
    real multi-distance, multi-angle capture remains for dedicated evaluation.
    \item \emph{Pose-extraction failure.} Under heavy occlusion MediaPipe may lose
    tracking, drifting predictions toward the cohort mean; confidence-based frame
    filtering is recommended for uncontrolled environments.
    \item \emph{Statistical testing.} Sequence-level significance tests inflate
    the effective sample size because sequences cluster within subjects;
    confirmatory analysis aggregates errors at the subject level
    (Section~\ref{sec:metrics}).
\end{itemize}

%% =======================================================
%%  7. CONCLUSION
%% =======================================================
\section{Conclusion}
\label{sec:conclusion}
We presented an intelligent vision-based decision-support system for Parkinson's
gait severity assessment from a single RGB camera. Beyond a bare motion encoder,
the system addresses practical deployment bottlenecks: it uses COM-centred
normalisation, validates an optional scale-normalised robustness variant under a
global-scaling approximation, enforces clinical output bounds, and exposes an
explicit abstention rule quantified by a risk--coverage analysis. Validated on the
multi-site CARE-PD subset and an independent prospective CNUH cohort, the system
achieves competitive accuracy (MAE\,=\,0.358, RMSE\,=\,0.564,
MedAE\,=\,0.147) and Macro-F1 against
state-of-the-art encoders at a fraction of their parameter count, with stable
zero-shot cross-site behaviour. By reporting edge-inference latency, Macro-F1
under class imbalance, and transparent failure modes, the pipeline forms a
practical, privacy-preserving building block for clinical decision support in
Parkinson's disease.

%% =======================================================
%%  CRediT
%% =======================================================
\section*{CRediT authorship contribution statement}
\textbf{Bokyung Kim}: Conceptualisation, Methodology, Software, Formal analysis,
Data curation, Writing -- original draft, Visualisation.
\textbf{Hieyong Jeong}: Conceptualisation, Supervision, Writing -- review and
editing, Funding acquisition.
\textbf{Md Ilias Bappi}: Investigation, Writing -- review and editing.
\textbf{Kyungbaek Kim}: Resources, Writing -- review and editing.

\section*{Declaration of competing interest}
The authors declare no known competing financial interests or personal
relationships that could have influenced the work reported in this paper.

\section*{Data availability}
Source code is available at
\url{https://github.com/bokyung08/parkinson_COM}. The CNUH clinical data cannot
be redistributed due to patient-privacy regulations; requests may be directed to
the corresponding author. CARE-PD is available from \cite{adeli2025carepd}.

\section*{Ethical approval}
This study was approved by the IRB of Chonnam National University Hospital
(CNUH-2025-203). Written informed consent was obtained from all participants, in
accordance with the 1964 Helsinki Declaration and its later amendments.

\section*{Acknowledgements}
This work was supported by the National Research Foundation of Korea (NRF) grant
funded by the Korea government (MSIT) (No.~RS-2026-25469002), and by the
IITP-ITRC grant funded by the Korea government (MSIT)
(IITP-2026-RS-2024-00437718).

%% =======================================================
%%  REFERENCES
%% =======================================================
\bibliographystyle{elsarticle-num}
\begin{thebibliography}{99}
\bibitem{martinez2011mds} Mart\'inez-Mart\'in P, et al. MDS-UPDRS: a new scale for the evaluation of Parkinson's disease. \textit{Mov Disord} 2011;26(12):2171--2177.
\bibitem{richards1994inter} Richards M, et al. Inter-rater reliability of the Unified Parkinson's Disease Rating Scale. \textit{Mov Disord} 1994;9(1):89--91.
\bibitem{goetz2008mds} Goetz CG, et al. MDS-UPDRS: scale presentation and clinimetric testing results. \textit{Mov Disord} 2008;23(15):2129--2170.
\bibitem{dibiase2020gait} di Biase L, et al. Gait analysis in Parkinson's disease: most accurate markers for diagnosis and monitoring. \textit{Sensors} 2020;20(12):3529.
\bibitem{zhang2024wearable} Zhang W, et al. Wearable sensor-based quantitative gait analysis in Parkinson's disease. \textit{npj Digit Med} 2024;7:169.
\bibitem{sun2024digital} Sun Y, et al. Digital biomarkers for precision diagnosis and monitoring in Parkinson's disease. \textit{npj Digit Med} 2024;7:218.
\bibitem{cao2017openpose} Cao Z, et al. OpenPose: realtime multi-person 2D pose estimation using part affinity fields. \textit{IEEE TPAMI} 2019;43(1):172--186.
\bibitem{bazarevsky2020blazepose} Bazarevsky V, et al. BlazePose: on-device real-time body pose tracking. \textit{arXiv} 2020;arXiv:2006.10204.
\bibitem{lu2020vision} Lu M, et al. Vision-based estimation of MDS-UPDRS gait scores for assessing Parkinson's disease motor severity. In: \textit{MICCAI}; 2020. p.\ 637--647.
\bibitem{lu2021quantifying} Lu M, et al. Quantifying Parkinson's disease motor severity under uncertainty using MDS-UPDRS videos. \textit{Med Image Anal} 2021;73:102179.
\bibitem{gennarelli2025skelmamba} Gennarelli I, et al. SkelMamba: a markerless motion capture approach for automated diagnosis of neurological disorders based on RGB videos. \textit{Gait Posture} 2025;122:109944.
\bibitem{tian2024crossspatiotemporal} Tian H, et al. Cross-spatiotemporal graph convolution networks for skeleton-based Parkinsonian gait MDS-UPDRS score estimation. \textit{IEEE TNSRE} 2024;32:412--421.
\bibitem{adeli2025carepd} Adeli V, et al. CARE-PD: a multi-site anonymized clinical dataset for Parkinson's disease gait assessment. In: \textit{NeurIPS Datasets and Benchmarks}; 2025.
\bibitem{adeli2024benchmarking} Adeli V, et al. Benchmarking skeleton-based motion encoder models for clinical applications: estimating Parkinson's disease severity in walking sequences. In: \textit{IEEE FG}; 2024.
\bibitem{yan2018spatial} Yan S, Xiong Y, Lin D. Spatial temporal graph convolutional networks for skeleton-based action recognition. In: \textit{AAAI}; 2018.
\bibitem{shi2019twoStream} Shi L, et al. Two-stream adaptive graph convolutional networks for skeleton-based action recognition. In: \textit{CVPR}; 2019.
\bibitem{zhu2023motionbert} Zhu W, et al. MotionBERT: a unified perspective on learning human motion representations. In: \textit{ICCV}; 2023.
\bibitem{mehraban2024motionagformer} Mehraban S, et al. MotionAGFormer: enhancing 3D human pose estimation with a transformer-GCNFormer network. In: \textit{WACV}; 2024.
\bibitem{winter2009biomechanics} Winter DA. \textit{Biomechanics and Motor Control of Human Movement}. 4th ed.\ Wiley; 2009.
\end{thebibliography}

\end{document}
